\begin{mistake}{2026-04-12}{06}

\begin{problemcontent}
设四阶行列式
\[
f(x)=\begin{vmatrix}
1&2&3&4\\
2&x&4&1\\
3&4&x&2\\
4&1&2&3
\end{vmatrix},
\]
求展开式中 \(x^2\) 项的系数。

\end{problemcontent}

\begin{solutioncontent}
矩阵中含 \(x\) 的元素只有 \(a_{22}=x\) 与 \(a_{33}=x\)。要得到 \(x^2\) 项，必须在行列式展开中同时选取这两个元素。

按广义 Laplace 展开，取第 2、3 行与第 2、3 列对应的子式：
\[
M=\begin{vmatrix}
x&4\\
4&x
\end{vmatrix}=x^2-16.
\]

删去第 2、3 行与第 2、3 列后，余下的余子式为
\[
N=\begin{vmatrix}
1&4\\
4&3
\end{vmatrix}=3-16=-13.
\]

符号因子为
\[
(-1)^{(2+3)+(2+3)}=(-1)^{10}=1.
\]

因此对应部分为
\[
f(x)\text{ 中由这两行两列产生的项}=1\cdot M\cdot N=(x^2-16)(-13).
\]

故 \(x^2\) 项的系数为
\[
-13.
\]

\end{solutioncontent}

\mistakeknowledge{
\begin{itemize}
  \item \textbf{核心公式/定理}：求某一特定高次项系数时，应只保留能产生该次数的含参元素组合。广义 Laplace 展开公式为“所取子式乘对应代数余子式”。
  \item \textbf{易错点}：余子式前必须乘符号因子 \((-1)^{\sum i_k+\sum j_k}\)。另一个高频错误是删行删列后把剩余元素的相对位置写错。
  \item \textbf{关联知识}：若参数只出现在少数几个位置，优先用“按参数元素定位取项”的方法求系数，而不要把整个四阶行列式完整展开。
\end{itemize}}

\end{mistake}
